\newpage

\section{量子力学：氢原子在脉冲电场中的电偶极跃迁}
\label{sec:hydrogen-pulse}

\begin{modelbox}[氢原子在脉冲电场中的电偶极跃迁]
基态氢原子在指数衰减电场中的跃迁。

基态氢原子处于弱电场中，微扰哈密顿量为
\[
\hat{H}' = \begin{cases}
0, & t \leq 0 \\[4pt]
\lambda z\, e^{-t/T}, & t > 0
\end{cases}
\]
其中 $\lambda, T$ 为常数。

\textbf{(1)} 求很长时间后 ($t \gg T$) 电子跃迁到激发态 $\psi_{210}$ 的概率。

\textbf{(2)} 基态电子跃迁到下列哪个激发态的概率等于零？简述理由。
\begin{itemize}
\item[A.] $\psi_{200}$
\item[B.] $\psi_{211}$
\item[C.] $\psi_{21,-1}$
\item[D.] $\psi_{210}$
\end{itemize}
\end{modelbox}

\subsection{考点快速分析}

\begin{tipbox}[考点快速分析]
\begin{itemize}
\item 含时微扰一级公式：$a_{fi}(t) = \dfrac{1}{i\hbar} \int_0^t \langle f|H'(t')|i\rangle e^{i\omega_{fi}t'} dt'$
\item 指数时间积分：$\int_0^\infty e^{-t/T} e^{i\omega t} dt = \dfrac{1}{1/T - i\omega} = \dfrac{T}{1 - i\omega T}$
\item 选择定则：$\Delta l = \pm 1$，$\Delta m = 0$ ($z$ 方向电场)
\item 角向积分：$Y_{00}$、$Y_{10}$、$\cos\theta \propto Y_{10}$ 的耦合
\end{itemize}
\end{tipbox}

\subsection{选择定则分析——第 (2) 问}

\begin{derivationbox}[选择定则分析]
微扰算符 $H' \propto z = r\cos\theta$ 是电偶极算符。在球谐函数表象中：
\begin{itemize}
\item $\cos\theta \propto Y_{10}$，角量子数为 1，磁量子数为 0
\item 基态 $|100\rangle$：$l=0$，$m=0$
\end{itemize}
根据选择定则：
\begin{enumerate}
\item $\Delta l = \pm 1$：末态必须有 $l=1$。选项 A ($\psi_{200}$，$l=0$) 不满足，概率为零。
\item $\Delta m = 0$：因为 $z$ 方向微扰不依赖于 $\phi$（即 $L_z$ 守恒）。选项 B ($\psi_{211}$，$m=1$) 和 C ($\psi_{21,-1}$，$m=-1$) 不满足，概率为零。
\end{enumerate}
因此只有 D ($\psi_{210}$，$l=1, m=0$) 是允许跃迁的通道。答案为 ABC。
\end{derivationbox}

\subsection{跃迁到 $\psi_{210}$ 的概率计算——第 (1) 问}

\begin{derivationbox}[跃迁概率计算]
\textbf{Step 1：矩阵元} $z_{21,10} = \langle 210|z|100\rangle$

波函数：
\[
\psi_{100} = \frac{2}{a^{3/2}} e^{-r/a} \frac{1}{\sqrt{4\pi}}, \qquad
\psi_{210} = \frac{1}{(2a)^{3/2}} \frac{r}{\sqrt{3}a} e^{-r/2a} \sqrt{\frac{3}{4\pi}} \cos\theta.
\]
角向积分（利用 $\cos\theta = \sqrt{\dfrac{4\pi}{3}} Y_{10}$）：
\[
\int Y_{10}^* \cos\theta \, Y_{00} \, d\Omega = \sqrt{\frac{4\pi}{3}} \frac{1}{\sqrt{4\pi}} \int |Y_{10}|^2 d\Omega = \frac{1}{\sqrt{3}}.
\]
径向积分：
\[
\int_0^\infty R_{21}(r)\, r\, R_{10}(r)\, r^2 dr = \frac{1}{\sqrt{6}\,a^4} \int_0^\infty r^4 e^{-3r/2a} dr = \frac{1}{\sqrt{6}\,a^4} \cdot \frac{4!}{(3/2a)^5} = \frac{256a}{81\sqrt{6}}.
\]
因此：
\[
z_{21,10} = \frac{256a}{81\sqrt{6}} \cdot \frac{1}{\sqrt{3}} = \frac{256a}{243\sqrt{2}} = \frac{128\sqrt{2}\,a}{243}.
\]

\textbf{Step 2：时间积分}

$H'_{21,10}(t) = \lambda z_{21,10}\, e^{-t/T}$，$\omega_{21} = (E_2 - E_1)/\hbar = -\dfrac{3e^2}{8a\hbar}$（注意符号，取绝对值计算概率）。
\[
a_{fi}(\infty) = \frac{\lambda z_{21,10}}{i\hbar} \int_0^\infty e^{-t/T} e^{i\omega_{21}t} dt = \frac{\lambda z_{21,10}}{i\hbar} \cdot \frac{T}{1 - i\omega_{21}T}.
\]

\textbf{Step 3：概率}
\[
P = |a_{fi}|^2 = \frac{\lambda^2 |z_{21,10}|^2}{\hbar^2} \cdot \frac{T^2}{1 + \omega_{21}^2 T^2}.
\]
代入 $|z_{21,10}|^2 = \dfrac{2^{15} a^2}{3^{10}}$：
\[
P = 2^5 \left(\frac{2}{3}\right)^{10} \frac{\lambda^2 a^2}{\hbar^2} \cdot \frac{T^2}{1 + \omega_{21}^2 T^2}.
\]
\end{derivationbox}

\subsection{径向积分的详细计算}

\begin{derivationbox}[径向积分的详细计算]
标准氢原子径向函数：
\[
R_{10}(r) = \frac{2}{a^{3/2}} e^{-r/a}, \qquad R_{21}(r) = \frac{1}{2\sqrt{6}\,a^{3/2}} \frac{r}{a} e^{-r/2a}.
\]
径向矩阵元：
\begin{align*}
\int_0^\infty R_{21}(r)\, r\, R_{10}(r)\, r^2 dr
&= \int_0^\infty \frac{1}{2\sqrt{6}\,a^{3/2}} \frac{r}{a} e^{-r/2a} \cdot r \cdot \frac{2}{a^{3/2}} e^{-r/a} \cdot r^2 dr \\
&= \frac{1}{\sqrt{6}\,a^4} \int_0^\infty r^4 e^{-3r/2a} dr.
\end{align*}
利用标准积分 $\int_0^\infty x^n e^{-\alpha x} dx = n!/\alpha^{n+1}$ ($n=4$，$\alpha = 3/2a$)：
\[
\int_0^\infty r^4 e^{-3r/2a} dr = \frac{24}{(3/2a)^5} = \frac{24 \cdot 32 a^5}{243} = \frac{768 a^5}{243} = \frac{256 a^5}{81}.
\]
因此径向积分为 $\dfrac{256a}{81\sqrt{6}}$。乘以角向因子 $1/\sqrt{3}$ 即得 $z_{21,10}$。
\end{derivationbox}

\subsection{结果的物理极限}

\begin{insightbox}[结果的物理极限]
\begin{itemize}
\item $\omega_{21} T \ll 1$（慢开关，$T$ 很短）：$P \approx \dfrac{\lambda^2 |z_{21,10}|^2 T^2}{\hbar^2}$，跃迁概率与作用时间的平方成正比。
\item $\omega_{21} T \gg 1$（快开关，$T$ 很长）：$P \approx \dfrac{\lambda^2 |z_{21,10}|^2}{\hbar^2 \omega_{21}^2}$，概率饱和，与 $T$ 无关。这是能量--时间不确定关系的表现：作用时间远大于内禀周期时，系统有时间调整，跃迁被抑制。
\end{itemize}
\end{insightbox}
